% ========================================================================================
% CAPITOLO 8 - PROVA D'ESAME DEL 19 DICEMBRE 2024
% ========================================================================================
\chapter{Prova d'esame del 19 dicembre 2024}

% ----------------------------------------------------------------------------------------
\section{Esercizio 12}
% ----------------------------------------------------------------------------------------

\begin{traccia}
	Un sistema è composto da 2 nodi, A e B. A include una ROM (progettata come macchina
	sequenziale con READ sincrono) di 8 locazioni da 4 bit, mentre B include un sommatore
	parallelo in grado di effettuare la somma di 2 stringhe di 4 bit ciascuna e un registro R
	di 4 bit. Il sistema opera come segue: all'arrivo di un segnale di start, A inizia a
	prelevare gli elementi ROM[i] dalla propria memoria e li invia, uno alla volta, a B
	mediante handshaking. B somma progressivamente le stringhe ricevute utilizzando il
	sommatore e alla fine inserisce il risultato nel registro R.
	\begin{enumerate}[nosep]
		\item Si disegni l'architettura complessiva del sistema tramite un diagramma a blocchi,
		      identificando parte operativa e parte di controllo di ciascun nodo. Ogni nodo deve
		      essere progettato seguendo un approccio strutturale, individuando tutti i
		      componenti, le loro interfacce e le loro interconnessioni.
		\item Si progettino le unità di controllo di A e B evidenziando gli stati, gli ingressi
		      e le uscite negli automi risultanti. È obbligatorio specificare la tempificazione
		      che si intende dare alle macchine (fronte attivo del clock, tempificazione dei
		      segnali di READ/WRITE su registri e memorie).
		\item Si progetti il sommatore secondo un'architettura di tipo carry look ahead.
		\item Si fornisca l'implementazione in VHDL dell'intero sistema e si proceda alla
		      simulazione nel caso in cui il clock del sistema A e del sistema B siano diversi
		      (A più lento e A più veloce).
	\end{enumerate}
\end{traccia}

\subsection{Architettura del sistema}
\todo{Diagramma a blocchi con parte operativa e parte di controllo di A e di B.}

\begin{attenzione}[Overflow del registro R]
	La somma di 8 valori da 4 bit può arrivare a $8 \times 15 = 120$, che non è
	rappresentabile su 4 bit. \todo{Dichiarare come è trattato: somma modulo 16 ed eventuale flag di overflow.}
\end{attenzione}

\subsection{Unità di controllo}
\todo{Automi di A e B con tempificazione esplicita: fronte attivo, READ della ROM, WRITE di R.}

\subsection{Sommatore carry look-ahead}

Per ogni bit si definiscono i segnali di \emph{generate} e \emph{propagate}
\begin{equation}
	G_i = a_i\, b_i, \qquad P_i = a_i \oplus b_i,
\end{equation}
da cui tutti i riporti si calcolano in parallelo a partire da $c_0$:
\begin{align}
	c_1 &= G_0 + P_0 c_0 \\
	c_2 &= G_1 + P_1 G_0 + P_1 P_0 c_0 \\
	c_3 &= G_2 + P_2 G_1 + P_2 P_1 G_0 + P_2 P_1 P_0 c_0 \\
	c_4 &= G_3 + P_3 G_2 + P_3 P_2 G_1 + P_3 P_2 P_1 G_0 + P_3 P_2 P_1 P_0 c_0
\end{align}
e la somma è $s_i = P_i \oplus c_i$.

\todo{Schema a porte, codice strutturale, testbench esaustivo.}

\subsection{Implementazione e simulazione}
\todo{Codice del sistema; simulazioni con $T_A > T_B$ e $T_A < T_B$; sincronizzazione di \sig{req} e \sig{ack}.}

\screenshot[0.95\linewidth]{es12_sim_a_lento.png}{Simulazione con il clock di A più lento di quello di B.}{fig:es12-lento}

\screenshot[0.95\linewidth]{es12_sim_a_veloce.png}{Simulazione con il clock di A più veloce di quello di B.}{fig:es12-veloce}

\subsection{Timing analysis}
\todo{Confronto del percorso critico fra RCA e CLA.}

\begin{table}[H]
	\centering\small
	\caption{Confronto fra sommatore ripple-carry e carry look-ahead (valori da Vivado).}
	\label{tab:rca-cla}
	\begin{tabular}{@{}lcccc@{}}
		\toprule
		\textbf{Architettura} & \textbf{LUT} & \textbf{Livelli di logica} & \textbf{Ritardo [\si{\nano\second}]} & \textbf{WNS [\si{\nano\second}]} \\ \midrule
		Ripple-carry     & \todo{} & & & \\
		Carry look-ahead & & & & \\
		\bottomrule
	\end{tabular}
\end{table}
